% ═══════════════════════════════════════════════════════════════════════════
% SCAFFOLD DS-04: Beispiel-E-Mail (Punto final)
% Modell für Gruppe A: E-Mail an Austauschschüler
% ═══════════════════════════════════════════════════════════════════════════

\documentclass[11pt, a4paper]{article}

\usepackage[utf8]{inputenc}
\usepackage[T1]{fontenc}
\usepackage[ngerman]{babel}
\usepackage{helvet}
\renewcommand{\familydefault}{\sfdefault}

\usepackage[margin=1.5cm, top=1.2cm, bottom=1.5cm]{geometry}
\usepackage[table]{xcolor}
\usepackage{array}
\usepackage{tabularx}
\usepackage{tikz}
\usepackage{tcolorbox}
\usepackage{enumitem}
\usepackage{amssymb}
\tcbuselibrary{skins}

% Farben
\definecolor{spanisch}{HTML}{c41e3a}
\definecolor{gruppeA}{HTML}{2a7a4b}
\definecolor{hellgrau}{HTML}{f5f5f5}
\definecolor{struktur}{HTML}{1565c0}

\pagestyle{empty}

\begin{document}

% ═══════════════════════════════════════════════════════════════════════════
% TITEL
% ═══════════════════════════════════════════════════════════════════════════

\begin{center}
\begin{tikzpicture}
\fill[gruppeA] (0,0) rectangle (18,0.9);
\node[white, font=\Large\bfseries] at (9,0.45) {Beispiel-E-Mail: Mi barrio (Punto final)};
\end{tikzpicture}
\end{center}

\vspace{0.3cm}

% ═══════════════════════════════════════════════════════════════════════════
% AUFGABE
% ═══════════════════════════════════════════════════════════════════════════

\begin{tcolorbox}[
    colback=hellgrau,
    colframe=spanisch,
    title={\color{white}\bfseries Aufgabe (Punto final, S. 33)},
    fonttitle=\bfseries
]
\textbf{Schreibe eine E-Mail an deinen spanischen Austauschschüler Pablo/Paula.}\\
Beschreibe dein Stadtviertel:
\begin{itemize}[leftmargin=2em]
    \item Welche Orte gibt es? (hay)
    \item Wo liegen sie? (está/están)
    \item Wie ist dein Viertel? (es)
    \item Was kann man dort machen? (poder, ir a)
\end{itemize}
Länge: ca. 80-100 Wörter
\end{tcolorbox}

\vspace{0.3cm}

% ═══════════════════════════════════════════════════════════════════════════
% STRUKTUR
% ═══════════════════════════════════════════════════════════════════════════

\begin{tcolorbox}[
    colback=struktur!10,
    colframe=struktur,
    title={\color{white}\bfseries Struktur einer E-Mail},
    fonttitle=\bfseries
]

\begin{tabular}{@{}p{4cm}p{12.5cm}@{}}
\textbf{1. Anrede} & ¡Hola Pablo! / Querido Pablo, / Querida Paula, \\[0.3cm]
\textbf{2. Einleitung} & ¿Qué tal? / ¿Cómo estás? / Espero que estés bien. \\[0.3cm]
\textbf{3. Hauptteil} & Beschreibung des Stadtviertels (2-3 Absätze) \\[0.3cm]
\textbf{4. Abschluss} & ¿Y tú? ¿Cómo es tu barrio? / Escríbeme pronto. \\[0.3cm]
\textbf{5. Gruß} & Un abrazo, / Hasta pronto, / Saludos, + \textit{[dein Name]} \\
\end{tabular}

\end{tcolorbox}

\vspace{0.3cm}

% ═══════════════════════════════════════════════════════════════════════════
% BEISPIEL
% ═══════════════════════════════════════════════════════════════════════════

\begin{tcolorbox}[
    colback=gruppeA!10,
    colframe=gruppeA,
    title={\color{white}\bfseries Beispiel-E-Mail (Muster)},
    fonttitle=\bfseries
]

\textbf{De:} maria@email.de\\
\textbf{Para:} pablo@email.es\\
\textbf{Asunto:} Mi barrio en Rostock

\vspace{0.3cm}

\textit{¡Hola Pablo!}

\textit{¿Qué tal? Espero que estés bien. Hoy quiero \underline{describir} mi barrio.}

\textit{Vivo en un barrio que \textbf{es} muy \underline{tranquilo} y \underline{bonito}. \textbf{Hay} muchas cosas aquí: \textbf{hay} un parque grande, un supermercado y una panadería. También \textbf{hay} una biblioteca muy moderna.}

\textit{El parque \textbf{está} cerca de mi casa, al lado del colegio. El supermercado \textbf{está} enfrente de la parada de autobús. La biblioteca \textbf{está} en el centro, entre el banco y la farmacia.}

\textit{Mi barrio \textbf{es} perfecto porque \textbf{puedo} ir a pie a todo. Los fines de semana \textbf{voy} al parque con mis amigos. También \textbf{podemos} ir al centro comercial en autobús.}

\textit{¿Y tú? ¿Cómo \textbf{es} tu barrio? ¿Qué \textbf{hay} cerca de tu casa?}

\textit{¡Escríbeme pronto!}

\textit{Un abrazo,}\\
\textit{María}

\vspace{0.2cm}
\footnotesize\textcolor{spanisch}{Wörter: ca. 120 | Verwendete Strukturen: hay (5x), es (4x), está (4x), poder (2x), ir (2x)}
\end{tcolorbox}

\vspace{0.3cm}

% ═══════════════════════════════════════════════════════════════════════════
% TIPPS
% ═══════════════════════════════════════════════════════════════════════════

\begin{center}
\fcolorbox{spanisch}{hellgrau}{%
\begin{minipage}{0.95\textwidth}
\textbf{Checkliste vor dem Abgeben:}
\begin{itemize}[leftmargin=2em]
    \item[$\square$] Anrede und Gruß vorhanden?
    \item[$\square$] Mind. 3x \textbf{hay} verwendet?
    \item[$\square$] Mind. 3x \textbf{estar} (Ortsangaben) verwendet?
    \item[$\square$] Mind. 2x \textbf{ser} (Beschreibungen) verwendet?
    \item[$\square$] Frage an den Austauschschüler gestellt?
    \item[$\square$] 80-100 Wörter geschrieben?
\end{itemize}
\end{minipage}%
}
\end{center}

\end{document}
